\documentclass[12pt,letterpaper, onecolumn]{exam}
\usepackage{amsmath}
\usepackage{amssymb}
\usepackage[lmargin=71pt, tmargin=1.2in]{geometry}  %For centering solution box
\usepackage{xcolor}
\usepackage[brazil]{babel}
\usepackage{natbib}   % very common in economics
\usepackage{enumitem}
\usepackage{booktabs}

\usepackage[hidelinks]{hyperref}
\urlstyle{same}



\renewenvironment{solution}
  {\par\noindent\textbf{R:}\enspace\color{blue}}
  {\par\color{black}}

% \chead{\hline} % Un-comment to draw line below header
\thispagestyle{empty}   %For removing header/footer from page 1

\begin{document}

\begingroup  
    \centering
    \LARGE Lista 2 Gabarito - Econometria I - 2026\\[0.5em]
    %\large \today\\[0.5em]
    \large Prof: André Luiz Pereira Mancha \par
    \large Monitor: Tuffy Licciardi Issa  \par
   \small \href{mailto:tuffy@usp.br}{tuffy@usp.br} \quad
\href{https://github.com/Tuffyli}{github.com/tuffyli} \par
    %\large Roll Number\par
    %\large Class/Batch/Section\par
\endgroup
\rule{\textwidth}{0.4pt}
\pointsdroppedatright   %Self-explanatory
\printanswers
\renewcommand{\solutiontitle}{\noindent\textbf{Ans:}\enspace}   %Replace "Ans:" with starting keyword in solution box


\begin{questions}

    \question (2.4 \citep{wooldridge2016}) Os dados do arquivo BWGHT contêm infomrações de nascimentos por mulheres nos Estados Unidos. As duas variáveis de interesse são: a variável dependente, peso dos recém-nascidos em onças (\textit{bwght}), e a variável explicativa, número de médio de cigarros que a mãe fumou por dia durante a gravidez (\textit{cigs}). A seguinte regressão simples foi estimada usando dados de $n = 1.388$ nascimentos.\begin{equation*}
        \widehat{bwght} = 119,77 - 0,514 \cdot cigs
    \end{equation*}
    \begin{enumerate}[label=(\alph*)]
    \item Qual é o peso de nascimento previsto quando $cigs = 0$? E quando $cigs = 20$ (um maço por dia)? Comente a diferença.
    \item O modelo de regressão simples necessariamente captura uma relação causal entre o peso de nascimento da criança e os hábitos de fumar da mãe? Explique.
    \item Para prever um peso de nascimento de 125 onças, qual deveria ser a magnitude de $cigs$? Comente.
    \item A proporção de mulheres que não fumaram durante a gravidez na amostra é de cerca de 0,85. Isso ajuda a reconciliar a sua conclusão do item (c)?
    \end{enumerate}

\begin{solution}
        \begin{enumerate}[label=(\alph*)]
    \item $\widehat{bwght}(0)=119{,}77$, \quad $\widehat{bwght}(20)=109{,}49$. Diferença de $10{,}28$ mais cigarros correlacionado com menor peso.
    \item Não. Lembrem-se da máxima \textbf{Correlação não implica causalidade}.
    \item  $cigs \approx -10{,}17$
    \item Sim. Com poucas observações fumantes, a regressão pode estar superestimando $\hat\beta_1 = 0,514$.
    \end{enumerate}
\end{solution}
    
    \question (2.8 \citep{wooldridge2016}) Considere o modelo de regressão simples padrão $y = \beta_0 + \beta_1x + u$ sob as hipóteses de Gauss-Markov RLS.1 a RLS.5. Os estimadores usuais $\hat{\beta}_0$ e $\hat{\beta}_1$ são não viesados para seus respectivos parâmetros populacionais. Seja $\tilde{\beta}_1$ o estimador de $\beta_1$ obtido ao assumir que o intercepto é zero.
    \begin{enumerate}[label=(\alph*)]
    \item Encontre E($\tilde{\beta}_1$) em termos de $x_i$, $\beta_0$ e $\beta_1$. Verifique que se $\tilde \beta_1$ é não viesado para $\beta_1$ quando o intercepto populacional ($\beta_0$) é zero. Há outros casos em que $\tilde \beta_1$ é não viesado?

    \item Encontre a variância de $\tilde \beta_1$. [Dica: a variância não depende de $\beta_0$.]

    \item Mostre que Var($\tilde \beta_1$) $\leq$ Var($\hat \beta_1$). [Dica: para qualquer amostra de dados, $\sum^n_{i = 1}x^2_i \geq \sum^n_{i=1}(x - \bar{x})^2$, com a desigualdade estrita preponderando, a não ser que $\bar{x} = 0$.]

    \item Comente a relação entre viés e variância ao escolher entre $\hat \beta_1$ e $\tilde \beta_1$.
    \end{enumerate}
\begin{solution}
    \begin{enumerate}[label = (\alph*)]
        \item $E[\tilde \beta_1] = \beta_0 \frac{\sum x_i}{\sum x_i^2} + \beta_1$. Em caso de $\beta_0 = 0$, temos que $E[\tilde \beta_1] = \beta_1$, ou seja, sem viés. Outro caso seria se $\sum x_i = 0$.


        
        \item Vamos fazer o passo a passo da derivação. Para obter a variância condicional em \(X\), note que os dois primeiros termos são constantes. Portanto,
            \[
            \mathrm{Var}(\tilde{\beta}_1 \mid X)
            =
            \mathrm{Var}\left(
            \frac{\sum_{i=1}^n x_i u_i}{\sum_{i=1}^n x_i^2}
            \;\middle|\;
            X
            \right).
            \]
            
            Como \(\sum_{i=1}^n x_i^2\) é constante condicionalmente a \(X\), podemos fatorá-lo:
            \[
            \mathrm{Var}(\tilde{\beta}_1 \mid X)
            =
            \frac{1}{\left(\sum_{i=1}^n x_i^2\right)^2}
            \mathrm{Var}\left(
            \sum_{i=1}^n x_i u_i
            \;\middle|\;
            X
            \right).
            \]
            
            Agora, como se trata da variância de uma soma ao longo das observações,
            \[
            \mathrm{Var}\left(
            \sum_{i=1}^n x_i u_i
            \;\middle|\;
            X
            \right)
            =
            \sum_{i=1}^n x_i^2 \mathrm{Var}(u_i\mid X)
            +
            2\sum_{i<j} x_i x_j \mathrm{Cov}(u_i,u_j\mid X).
            \]
            
            Sob as hipóteses usuais do modelo,
            \[
            \mathrm{Var}(u_i\mid X)=\sigma^2
            \quad \text{e} \quad
            \mathrm{Cov}(u_i,u_j\mid X)=0 \ \text{para } i\neq j.
            \]
            
            Logo,
            \[
            \mathrm{Var}\left(
            \sum_{i=1}^n x_i u_i
            \;\middle|\;
            X
            \right)
            =
            \sigma^2 \sum_{i=1}^n x_i^2.
            \]
            
            Substituindo de volta,
            \[
            \mathrm{Var}(\tilde{\beta}_1 \mid X)
            =
            \frac{1}{\left(\sum_{i=1}^n x_i^2\right)^2}
            \left(
            \sigma^2 \sum_{i=1}^n x_i^2
            \right)
            =
            \frac{\sigma^2}{\sum_{i=1}^n x_i^2}.
            \]
            
            Portanto,
            \[
            \boxed{
            \mathrm{Var}(\tilde{\beta}_1 \mid X)
            =
            \frac{\sigma^2}{\sum_{i=1}^n x_i^2}
            }.
            \]
        \item Vale que
            \[
            \sum_{i=1}^n x_i^2
            =
            \sum_{i=1}^n (x_i-\bar{x})^2 + n\bar{x}^2.
            \]
            
            Como \(n\bar{x}^2 \geq 0\), segue que
            \[
            \sum_{i=1}^n x_i^2 \geq \sum_{i=1}^n (x_i-\bar{x})^2.
            \]
            Onde $\sum^n_{i = 1}x^2_i \geq \sum^n_{i=1}(x - \bar{x})^2$, pode ser transformado em:\[\frac{1}{\sum^n_{i = 1}x^2_i} \leq \frac{1}{\sum^n_{i=1}(x - \bar{x})^2}\]
            de modo que $Var(\tilde{\beta}_1) =\frac{\sigma^2}{\sum x_i^2} \le \frac{\sigma^2}{\sum (x_i-\bar x)^2} = Var(\hat{\beta}_1)$ \hfill $\blacksquare$
        \item Há um trade-off entre viés e variância. $\tilde \beta_1$ pode ser mais preciso, mas se $\beta_0 \neq 0$ e $\sum x_i \neq 0$ há viés. Sem a devida justificativa é arriscado impor o intercepto igual a zero na sua regressão.
        
    \end{enumerate}
\end{solution}



    \question (2.2 \citep{wooldridge2010}) A tabela seguinte contém variáveis $supGPA$ (nota média em curso superior nos Estados Unidos) e \textit{ACT} (nota do teste de avaliação de conhecimentos para ingresso em curso superior nos Estados Unidos) com as notas hipotéticas de oito estudantes de curso superior. O \textit{GPA} está baseado em uma escala de quatro pontos e foi arredondado para um dígito após o ponto decimal. A nota \textit{ACT} baseia-se em uma escala de 36 pontos e foi arredondada para um número inteiro.

    \begin{table}[bhtp]
        \centering
        \begin{tabular}{c|cc}
            \toprule
            \midrule
            \multicolumn{1}{c}{Estudante} & $supGPA$  & \textit{ACT} \\
            \midrule
             1 & 2,8 & 21\\
             2 & 3,4 & 24\\
             3 & 3,0 & 26\\
             4 & 3,5 & 27\\
             5 & 3,6 & 29\\
             6 & 3,0 & 25\\
             7 & 2,7 & 25\\
             8 & 3,7 & 30\\
             \midrule
             \bottomrule
        \end{tabular}
    \end{table}
    \begin{enumerate} [label = (\alph*)]
        \item Estime a relação entre \textit{GPA} e \textit{ACT} usando MQO; isto é, obtenha as estimativas de intercepto e de inclinação da equação:\begin{equation*}
            \widehat{GPA} = \hat \beta_0 + \hat \beta_1 act
        \end{equation*} Comente a direção da relação. O intercepto tem uma interpretação útil aqui? Explique. Qual deveria ser o valor previsto pelo \textit{GPA} se a nota \textit{ACT} aumentasse em cinco pontos?
        \item Calcule os valores estimados e os resíduos de cada observação e verifique que a soma dos resíduos é (aproximadamente) zero.
        \item Qual o valor previsto de ${GPA}$ quando $ACT = 20$?
        \item Quanto da variação de \textit{GPA} dos 8 estudantes é explicado pelo \textit{ACT}? Explique.
    \end{enumerate}
\begin{solution}
    \begin{enumerate}[label = (\alph*)]
        \item $\hat\beta_1 = 0,1022$ ,$\hat\beta_o = 0,5681$. O intercepto exibe o GPA caso um aluno tire 0 em \textit{ACT}, um cenário improvável dada a necessidade desta avaliação para o ingresso nos cursos e pela a amostra já ser composta de alunos do ensino superior. Direção positiva: maior nota no \textit{ACT} relacionada com maior \textit{GPA}. Aumento de 5 pontos seria $0,511$ no GPA.
        \item 0,0002 $\approx 0$.

    \begin{table}[h]
            \centering
            \caption{Resíduos}
            \begin{tabular}{c|c}
            \toprule
            \midrule
               1  & 0,0857 \\
               2  &  0,3791\\
               3  &  0,2253\\
               4  &  0,1725\\
               5  &  0,0681\\
               6  &  -0,1231\\
               7  &  -0,4231\\
               8  &  0,0659\\
               \midrule
               \bottomrule
            \end{tabular}
            \label{tab:placeholder}
        \end{table}
\item $\widehat{GPA}(20)=2{,}6121$.
\item $R^2 \approx 0{,}577$. Modos de calcular o $R^2$:\begin{equation}
    R^2 = \frac{SQE}{SQT} = \frac{\sum^n_{i=1}(\hat y_i - \bar y)^2}{\sum(y_i - \bar y)^2}
\end{equation}
\begin{equation}
    R^2 = 1- \frac{SQR}{SQT} =1- \frac{\sum^n_{i=1}(\hat u_i )^2}{\sum^n_{i=1}(y_i - \bar y)^2}
\end{equation}
O R-quadrado é a razão entre a variação explicada e a variação total.
\end{enumerate}

\end{solution}

\question (12 (ANPEC, 2005)) Um pesquisador estima o seguinte modelo de regressão simples: $Y_i = \beta_0 + \beta_1X_i + e_i$. Outro pesquisador estima o mesmo modelo, mas com escalas diferentes para $Y_i$ e $X_i$. O segundo modelo é: $Y^*_i = \beta_o^* + \beta_1^*X^*_i + e_i^*$, em que $Y^*_i = \omega_1Y_i$, $X^*_i = \omega_2X_i$, e $\omega_1$ e $\omega_2$ são constantes maiores do que zero. Responda com F ou V e explique.
\begin{enumerate}[label=(\arabic*), start=0]
    \item Os estimadores de Mínimos Quadrados Ordinários de $\beta_0$ e $\beta_1$ são iguais aos de $\beta_0^*$ e $\beta_1^*$
    \item Se $\hat \sigma^{*2}$ é a variância estimada de $e^*_i$ e $\hat \sigma^2$ é a variância estimada de $e_i$, então $\hat\sigma^{*2} = w^2_1\hat\sigma^2$.
    \item As variâncias dos estimadores dos parâmetros no primeiro modelo são maiores do que as variâncias dos estimadores do segundo modelo.
    \item Os coeficientes de determinação são iguais nos dois modelos.
    \item A transformação de escalada de ($Y_i$, $X_i$) para ($Y_i^*$, $X_i^*$) não afeta as propriedades de Mínimos Quadrados Ordinários dos parâmetros.
\end{enumerate}
\begin{solution}
    \begin{enumerate}[label=(\arabic*), start = 0]
        \item F
        \item V
        \item F
        \item V
        \item V
    \end{enumerate}
\end{solution}

\question (2.4 \citep{wooldridge2010}) Suponha que você está interessado em estimar o efeito das horas gastas com um curso de preparação para o vestibular (\textit{horas}) no total das notas do vestibular (\textit{ENEM}). A população é composta por todos os pré-universitários graduados no ensino médio em determinado ano.
    \begin{enumerate} [label = (\alph*)]
\item Suponha que lhe tenha sido dada uma subvenção para executar um experimento controlado. Explique como você estruturaria o experimento de forma a estimar o efeito causal de \textit{horas} no ENEM.
\item Considere o caso mais realístico em que os alunos decidem quanto tempo gastarão com um curso de preparação, e você só pode fazer amostragens aleatórias do \textit{ENEM} e \textit{horas} da população. Escreva o modelo populacional da seguinte forma:\begin{equation*}
    ENEM = \beta_0 + \beta_1 horas + u
\end{equation*} em que, como sempre, num modelo com um intercepto, podemos assumir $E[u] = 0$. Liste pelo menos dois fatores contidos em $u$. Existe a probabilidade de eles terem correlação negativa ou positiva com \textit{horas}?
\item Na equação da parte (b), qual deveria ser o sinal de $\beta_1$ se o curso de preparação for eficaz?
\item Na equação da parte (b), qual é a interpretação de $\beta_0$?
    \end{enumerate}
\begin{solution}
    \begin{enumerate} [label = (\alph*)]
        \item Aleatorização dos alunos participantes do curso preparatório do vestibular (grupo tratado e controle comparáveis). Depois comparar as médias finais entre grupos.
        \item Dois fatores são, habilidade inata e renda familiar. Ambos fatores provavelmente apresentam correlação positiva.
        \item $\beta_1 > 0$
        \item $\beta_0$ é a nota média no ENEM para um aluno com $horas = 0$
        
    \end{enumerate}
\end{solution}
    
\end{questions}
\bibliographystyle{apalike}   % or plainnat, econ, etc.
\bibliography{refs}
\end{document}